\documentclass[10pt,a4paper]{amsart}
\usepackage[margin=19mm]{geometry}
\usepackage{amsmath,amssymb,mathtools}
\usepackage[scr=rsfs,cal=euler]{mathalfa}
\usepackage{xfrac}
\usepackage[unicode,hidelinks,bookmarks=false]{hyperref}
\newcommand{\ie}{i.e.\ }
\newcommand{\cf}{cf.\ }
\newcommand{\resp}{respectively\ }
\newcommand{\Subsection}{\subsection}
\providecommand{\nobreakdash}{\nobreak\hbox{-}}
\setlength{\emergencystretch}{2em}
\multlinegap0pt
\usepackage{amssymb}
\usepackage{xcolor}
\usepackage{enumerate}
\newtheorem{lem}{Lemma}
\newcommand \eps{\varepsilon}
\title{An improvement of regularity result for pseudo Calabi flow}
\author{Jingrui Cheng, Junhao Tian}
\date{Source: arXiv:2603.19405v1 (CC BY 4.0)}
\begin{document}
\maketitle

\noindent\emph{Source and licence.} An improvement of regularity result for pseudo Calabi flow. Version arXiv:2603.19405v1 (CC BY 4.0), distributed by arXiv under the Creative Commons Attribution 4.0 International licence, http://creativecommons.org/licenses/by/4.0/, as recorded in the arXiv licence metadata for this exact version and shown on the abstract page https://arxiv.org/abs/2603.19405v1. Full licence notice: \texttt{licenses/arxiv-2603.19405-CC-BY-4.0.txt}.\par\smallskip

\noindent\textit{Editorial scope.} Counted unit: the article's lem Lem315 (source0.tex lines 1357--1362). The article's complete original proof of it is delivered with it (source0.tex lines 1364--1553, one \texttt{proof} environment of 190 lines). No mathematics is written, repaired or re-derived: the statement and the proof are the article's own lines, in the article's own notation, and no retained line is substituted or re-flowed.  No further source-local context is needed: every cross-reference of every retained line resolves inside the two delivered spans.  Nothing was omitted from any retained span, so the package carries no omission record.  No retained line contains a citation, so the wrapper carries no bibliography and no bibliography entry.  No retained line contains a figure environment, an image-inclusion command or drawing code, so no image is delivered and the package declares no asset dependency.  Editorial formatting only, in addition: an AMS \texttt{amsart} preamble that restates the article's own declarations and macros used by the retained lines, namely the environments \texttt{lem} on separate counters, together with the standard packages those lines need.  The retained environments print in this document as Lemma 1 in order of appearance, whereas the article prints the same items with the numbering of its own full document.  The complete native source of the article as distributed in the arXiv e-print for this exact version is bundled unchanged as \texttt{sources/196768/source0.tex}.
\begin{lem}[\textbf{Estimate of $|\nabla_\varphi P|^2$}]\label{Lem315}
	Let $\varphi$ be a smooth $\omega_0$-psh function.  Let $P$ solve $\Delta_{\varphi}P=-tr_{\omega_{\varphi}}Ric(\omega_0)+\underline{R}$,  and $\int_MP\omega_{\varphi}^n=0$.  Let $p>1$,  $K\ge 2$ large enough,  there exists $\delta_*>0$ small enough,  depending on $p$,  background metric,  and $||\log\frac{\omega_{\varphi}^n}{\omega_0^n}||_0$ and the choice of $K$,  such that if $P$ is $\delta_*$-close to a continuous function with modulus $\omega(r)$,  one has 
	\begin{align}
		K\int_M|\nabla_{\varphi}P|^{2(p+1)}\omega_{\varphi}^n\le C(p)\big(\int_M(1+tr_{\omega_{\varphi}}\omega_0)^{p+1}\omega_{\varphi}^n+\int_M|\nabla_{\varphi}F|^{2(p+1)}\omega_{\varphi}^n\big)+C. 
	\end{align}\\
\end{lem}

\begin{proof}[Proof of Lemma \ref{Lem315}]
 Let $K>0$,  we compute:
\begin{align}
	& \Delta_{\varphi}\big(e^{K(P-\tilde P)^2}|\nabla_{\varphi}P|^2\big)e^{-K(P-\tilde P)^2} =\left( \Delta_{\varphi}(K(P-\tilde P)^2)+|\nabla_{\varphi}(K(P-\tilde P)^2)|^2 \right)(|\nabla_\varphi P|^2)  \notag \\
	&\qquad\qquad\qquad\qquad\qquad +\Delta_{\varphi}(|\nabla_{\varphi}P|^2) +2\nabla_{\varphi}(K(P-\tilde P)^2)\cdot \nabla_{\varphi}(|\nabla_{\varphi}P|^2).
\end{align}\\
In the above,  we have:
\begin{align}
\Delta_{\varphi}(|\nabla_{\varphi}P|^2)=g_{\varphi}^{i\bar{q}}g_{\varphi}^{p\bar{j}}P_{,iq}P_{,\bar{p}\bar{j}}+g_{\varphi}^{i\bar{q}}g_{\varphi}^{p\bar{j}}P_{i\bar{j}}P_{p\bar{q}}+2\nabla_{\varphi}P\cdot \nabla_{\varphi}\Delta_{\varphi}P+g_{\varphi}^{i\bar{q}}g_{\varphi}^{p\bar{j}}(Ric(\omega_{\varphi}))_{i\bar{j}}P_pP_{\bar{q}}.
\end{align}\\
Also
\begin{align}
2\nabla_{\varphi}(K(P-\tilde P)^2)\cdot \nabla_{\varphi}(|\nabla_{\varphi}P|^2)=2Re\big(g_{\varphi}^{i\bar{j}}g_{\varphi}^{p\bar{q}}((K(P-\tilde P)^2)_iP_{p\bar{j}}P_{\bar{q}}+(K(P-\tilde P)^2)_iP_pP_{,\bar{q}\bar{j}})\big).
\end{align}\\
We have the following completion of squares:
\begin{align}
|\nabla_{\varphi}(K(P-\tilde P)^2)|^2 |\nabla_\varphi P|^2 +2Re\big(g_{\varphi}^{i\bar{j}}g_{\varphi}^{p\bar{q}}(K(P-\tilde P)^2)_iP_pP_{,\bar{q}\bar{j}}\big)+g_{\varphi}^{i\bar{q}}g_{\varphi}^{p\bar{j}}P_{,ip}P_{,\bar{q}\bar{j}}\ge  0
\end{align}\\
And 
\begin{align}
	4|\nabla_{\varphi}(K(P-\tilde P)^2)|^2|\nabla_\varphi P|^2 + 2Re\big(g_{\varphi}^{i\bar{j}}g_{\varphi}^{p\bar{q}}((K(P-\tilde P)^2)_iP_{p\bar{j}}P_{\bar{q}}) +\frac{1}{4}g_{\varphi}^{i\bar{q}}g_{\varphi}^{p\bar{j}}P_{i\bar{j}}P_{p\bar{q}}\ge  0
\end{align}\\
So we obtain:
\begin{align}
	&\Delta_{\varphi}\big(e^{K(P-\tilde P)^2}|\nabla_{\varphi}P|^2\big)e^{-K(P-\tilde P)^2} \ge  (2K-16K^2(P-\tilde P)^2)|\nabla_\varphi(P-\tilde P)|^2 |\nabla_{\varphi}P|^2  \notag \\
	&\qquad\qquad\qquad\qquad +2K(P-\tilde P)\Delta_{\varphi}(P-\tilde P)|\nabla_{\varphi}P|^2 +\frac{3}{4}g_{\varphi}^{i\bar{q}}g_{\varphi}^{p\bar{j}}P_{i\bar{j}}P_{p\bar{q}} \notag\\
	&\qquad\qquad\qquad\qquad +2\nabla_{\varphi}P\cdot \nabla_{\varphi}\Delta_{\varphi}P+g_{\varphi}^{i\bar{q}}g_{\varphi}^{p\bar{j}}(Ric(\omega_{\varphi}))_{i\bar{j}}P_pP_{\bar{q}} 
\end{align}\\
In the above,  we can estimate: 
\begin{align}
	&(2K-16K^2(P-\tilde P)^2)|\nabla_\varphi(P-\tilde P)|^2 |\nabla_{\varphi}P|^2 \notag\\
	&\quad  \ge  \big(2K-16K^2(\delta_*+\omega(r))^2 \big)\bigg( |\nabla_\varphi P|^4 -\frac{C}{r^2}|\nabla_\varphi P|^2 \bigg)
\end{align}\\
And
\begin{align}
	2K(P-\tilde P)\Delta_{\varphi}(P-\tilde P)|\nabla_{\varphi}P|^2\ge  -K(\delta_*+\omega(r)) \frac{C}{r^2} (tr_{\omega_{\varphi}}\omega_0+1)|\nabla_\varphi P|^2
\end{align}\\
Also
\begin{align}
	& g_{\varphi}^{i\bar{q}}g_{\varphi}^{p\bar{j}}(Ric(\omega_{\varphi}))_{i\bar{j}}P_pP_{\bar{q}} =g_{\varphi}^{i\bar{q}}g_{\varphi}^{p\bar{j}}(Ric(\omega_0))_{i\bar{j}}P_pP_{\bar{q}} - g_{\varphi}^{i\bar{q}}g_{\varphi}^{p\bar{j}}F_{i\bar{j}}P_pP_{\bar{q}} \notag \\
	&\qquad \ge  -C tr_{\omega_{\varphi}}\omega_0|\nabla_{\varphi}P|^2 - \Delta_{\varphi}F|\nabla_{\varphi}P|^2 - \frac{dd^cF\wedge d^c P\wedge d P\wedge \omega_{\varphi}^{n-2}}{\omega_{\varphi}^n}.
\end{align}\\
So we get:
\begin{align}
	&\Delta_{\varphi}\big(e^{K(P-\tilde P)^2}|\nabla_{\varphi}P|^2\big)e^{-K(P-\tilde P)^2}  \notag\\
	&\quad  \ge  \big(2K-16K^2(\delta_*+\omega(r))^2 \big)\bigg( |\nabla_\varphi P|^4 -\frac{C}{r^2}|\nabla_\varphi P|^2 \bigg)   \notag \\
	&\qquad +\frac{3}{4}g_{\varphi}^{i\bar{q}}g_{\varphi}^{p\bar{j}}P_{i\bar{j}}P_{p\bar{q}} -C tr_{\omega_{\varphi}}\omega_0|\nabla_{\varphi}P|^2 -K(\delta_*+\omega(r)) \frac{C}{r^2} (tr_{\omega_{\varphi}}\omega_0+1)|\nabla_\varphi P|^2 \notag\\
	&\qquad +2\nabla_{\varphi}P\cdot \nabla_{\varphi}\Delta_{\varphi}P - \Delta_{\varphi}F|\nabla_{\varphi}P|^2 - \frac{dd^cF\wedge d^cP\wedge dP\wedge \omega_{\varphi}^{n-2}}{\omega_{\varphi}^n}.
\end{align}\\
Denote $u=e^{K(P-\tilde{P})^2}|\nabla_{\varphi}P|^2$,  then the above inequality is equivalent to:
\begin{align}\label{4.133New}
	&\Delta_{\varphi}u  \ge  (2K-16K^2(\delta_*+\omega(r))^2)\bigg( |\nabla_\varphi P|^2 - \frac{C}{r^2}\bigg) u    \notag \\
	&\quad +\frac{3}{4}e^{K(P-\tilde P)^2} g_{\varphi}^{i\bar{q}}g_{\varphi}^{p\bar{j}}P_{i\bar{j}}P_{p\bar{q}} -C tr_{\omega_{\varphi}}\omega_0 u -K(\delta_*+\omega(r)) \frac{C}{r^2} (tr_{\omega_{\varphi}}\omega_0+1) u \notag\\
	&\quad+2e^{K(P-\tilde P)^2}\nabla_{\varphi}P\cdot\nabla_{\varphi}\Delta_{\varphi}P - \Delta_{\varphi}Fu - e^{K(P-\tilde P)^2}\frac{dd^cF\wedge d^cP\wedge dP\wedge \omega_{\varphi}^{n-2}}{\omega_{\varphi}^n}.
\end{align}\\

\noindent
Let $p>1$,  we have:
\begin{align}
	\Delta_\varphi u^p = p u^{p-1}\Delta_\varphi u+p(p-1)u^{p-2}|\nabla_\varphi u|^2.
\end{align}

Integrate with respect to $\omega_\varphi^n$, we have
\begin{equation}\label{4.135New}
\begin{split}
&\int_Mp(p-1)u^{p-2}|\nabla_{\varphi}u|^2\omega_{\varphi}^n=-\int_Mpu^{p-1}\Delta_{\varphi}u\omega_{\varphi}^n\\
&\le (2K-16K^2(\delta_*+\omega(r))^2)\bigg(-\int_Mpu^p|\nabla_{\varphi}P|^2\omega_{\varphi}^n+\frac{C}{r^2}\int_Mpu^p\omega_{\varphi}^n\bigg)\\
&-\int_M\frac{3p}{4}u^{p-1}e^{K(P-\tilde{P})^2}|\sqrt{-1}\partial\bar{\partial}P|^2_{\varphi}\omega_{\varphi}^n+(K(\delta_*+\omega(r))+1)\frac{C}{r^2}\int_Mpu^p(tr_{\omega_{\varphi}}\omega_0+1)\omega_{\varphi}^n\\
&-\int_M2pu^{p-1}e^{K(P-\tilde{P})^2}\nabla_{\varphi}P\cdot \nabla_{\varphi}\Delta_{\varphi}P\omega_{\varphi}^n+\int_Mpu^p\Delta_{\varphi}F\omega_{\varphi}^n\\
&+\int_Mpu^{p-1}e^{K(P-\tilde{P})^2}dd^cF\wedge d^cP\wedge dP\wedge \omega_{\varphi}^{n-2}.
\end{split}
\end{equation}
We are going to integrate by parts in the last three terms of the above.
\begin{align}\label{4.135N}
	&-\int_M2pu^{p-1}e^{K(P-\tilde P)^2}\nabla_{\varphi}P\cdot \nabla_{\varphi}\Delta_{\varphi}P\omega_{\varphi}^n =\int_M2p(p-1)u^{p-2}e^{K(P-\tilde P)^2}\nabla_{\varphi}u\cdot \nabla_{\varphi}P\Delta_{\varphi}P\omega_{\varphi}^n \notag \\
	&\quad\qquad\qquad +\int_M2pu^{p-1}e^{K(P-\tilde P)^2}2K (P-\tilde P)\left( |\nabla_{\varphi}P|^2-\nabla_\varphi\tilde P\cdot\nabla_\varphi P \right) \Delta_{\varphi}P\omega_{\varphi}^n \notag\\
	&\qquad\qquad\qquad\qquad + \int_M2pu^{p-1}e^{K(P-\tilde P)^2}(\Delta_{\varphi}P)^2\omega_{\varphi}^n.
\end{align}\\
Also
\begin{equation}\label{4.136}
\int_Mpu^p\Delta_{\varphi}F\omega_{\varphi}^n=-\int_Mp^2u^{p-1}\nabla_{\varphi}u\cdot \nabla_{\varphi}F\omega_{\varphi}^n.
\end{equation}

Next:
\begin{equation}\label{4.137}
\begin{split}
&\int_Mpu^{p-1}e^{K(P-\tilde{P})^2}dd^cF\wedge d^cP\wedge dP\wedge \omega_{\varphi}^{n-2}\\
&=-\int_Mp(p-1)u^{p-2}e^{K(P-\tilde{P})^2}du\wedge d^cF\wedge d^cP\wedge dP\wedge \omega_{\varphi}^{n-2}\\
&-\int_Mpu^{p-1}e^{K(P-\tilde{P})^2}2K(P-\tilde{P})d(P-\tilde{P})\wedge d^cF\wedge d^cP\wedge dP\wedge \omega_{\varphi}^{n-2}\\
&-\int_Mpu^{p-1}e^{K(P-\tilde{P})^2}d^cF\wedge dd^cP\wedge dP\wedge \omega_{\varphi}^{n-2}.
\end{split}
\end{equation}

Now we wish to estimate the right hand sides of (\ref{4.135N})-(\ref{4.137}).
We start with the right hand side of (\ref{4.135N}):
\begin{align}
	&\int_M2p(p-1)u^{p-2}e^{K(P-\tilde P)^2}\nabla_{\varphi}u\cdot \nabla_{\varphi}P\Delta_{\varphi}P\omega_{\varphi}^n \notag \\
	&\le  \frac{p(p-1)}{16}\int_Mu^{p-2}|\nabla_{\varphi}u|^2\omega_{\varphi}^n + C \int_M p(p-1)u^{p-1}e^{K(P-\tilde P)^2}(1+tr_{\omega_{\varphi}}\omega_0)^2\omega_{\varphi}^n.
\end{align}\\
Next
\begin{align}
	\int_M2pu^{p-1}e^{K(P-\tilde P)^2}(\Delta_{\varphi}P)^2\omega_{\varphi}^n\le  C \int_M pu^{p-1}e^{K(P-\tilde P)^2}(1+tr_{\omega_{\varphi}}\omega_0)^2\omega_{\varphi}^n.
\end{align}
Next,  keeping in mind that $u=e^{K(P-\tilde{P})^2}|\nabla_{\varphi}P|^2$ and that $|P-\tilde{P}|\le \delta_*+\omega(r)$.
\begin{align}
	&\int_M pu^{p-1}e^{K(P-\tilde P)^2} (P-\tilde P)\left( |\nabla_{\varphi}P|^2-\nabla_\varphi\tilde P\cdot\nabla_\varphi P \right) \Delta_{\varphi}P\omega_{\varphi}^n \notag\\
	& \le  \int_M pu^p |P-\tilde P||\Delta_\varphi P|\omega_\varphi^n + \int_M pu^{p-1}e^{K(P-\tilde P)^2} |P-\tilde P|\left( |\nabla_{\varphi}P|^2 +|\nabla_\varphi\tilde P|^2 \right) |\Delta_{\varphi}P| \omega_{\varphi}^n \notag\\
	&  \le  C(\delta_*+\omega(r)) \int_M pu^p  (1+tr_{\omega_{\varphi}}\omega_0)\omega_\varphi^n + \big(\delta_*+\omega(r)\big)\frac{C}{r^2} \int_M pu^{p-1} e^{K(P-\tilde P)^2}(1+tr_{\omega_{\varphi}}\omega_0)\omega_\varphi^n 
\end{align}\\

In the above,  we noted that $|\nabla_{\varphi}\tilde{P}|^2\le tr_{\omega_{\varphi}}\omega_0|\nabla \tilde{P}|^2\le tr_{\omega_{\varphi}}\omega_0\cdot \frac{C}{r^2}$.

Summing up we have 
\begin{align}\label{4.142}
	& -\int_M2pu^{p-1}e^{K(P-\tilde P)^2}\nabla_{\varphi}P\cdot \nabla_{\varphi}\Delta_{\varphi}P\omega_{\varphi}^n \le  \frac{p(p-1)}{16}\int_Mu^{p-2}|\nabla_{\varphi}u|^2\omega_{\varphi}^n \notag\\
	&\quad\quad + C \int_M p^2u^{p-1}e^{K(P-\tilde P)^2}(1+tr_{\omega_{\varphi}}\omega_0)^2\omega_{\varphi}^n +CK(\delta_*+\omega(r)) \int_M pu^p  (1+tr_{\omega_{\varphi}}\omega_0)\omega_\varphi^n \notag\\
	&\quad\quad\quad\quad\quad\quad\quad + K(\delta_*+\omega(r))\frac{C}{r^2} \int_M pu^{p-1} e^{K(P-\tilde P)^2}(1+tr_{\omega_{\varphi}}\omega_0)\omega_\varphi^n 
\end{align}

Now we look at (\ref{4.136}) and we can estimate:
\begin{equation}\label{4.143}
-\int_Mp^2u^{p-1}\nabla_{\varphi}u\cdot \nabla_{\varphi}F\omega_{\varphi}^n\le \int_M\frac{p(p-1)}{16}u^{p-2}|\nabla_{\varphi}P|^2+\int_M\frac{4p^3}{p-1}u^p|\nabla_{\varphi}F|^2\omega_{\varphi}^n.
\end{equation}
Next we are going to estimate the right hand side of (\ref{4.137}).
We have:
\begin{equation*}
\begin{split}
&-\int_Mp(p-1)u^{p-2}e^{K(P-\tilde{P})^2}du\wedge d^cF\wedge d^cP\wedge dP\wedge \omega_{\varphi}^{n-2}\\
&\le \int_Mp(p-1)u^{p-2}e^{K(P-\tilde{P})^2}|\nabla_{\varphi}u||\nabla_{\varphi}F||\nabla_{\varphi}P|^2\omega_{\varphi}^n=\int_Mp(p-1)u^{p-1}|\nabla_{\varphi}u||\nabla_{\varphi}F|\omega_{\varphi}^n\\
&\le \int_M\frac{p(p-1)u^{p-2}}{16}|\nabla_{\varphi}u|^2\omega_{\varphi}^n+\int_M4p(p-1)u^p|\nabla_{\varphi}F|^2\omega_{\varphi}^n.
\end{split}
\end{equation*}
Then
\begin{equation*}
\begin{split}
&-\int_Mpu^{p-1}e^{K(P-\tilde{P})^2}2K(P-\tilde{P})d(P-\tilde{P})\wedge d^cF\wedge d^cP\wedge dP\wedge \omega_{\varphi}^{n-2}\\
&\le \int_Mpu^{p-1}e^{K(P-\tilde{P})^2}2K(\delta_*+\omega(r))(|\nabla_{\varphi}P|+|\nabla_{\varphi}\tilde{P}|)|\nabla_{\varphi}F||\nabla_{\varphi}P|^2\omega_{\varphi}^n\\
&\le\int_Mpu^p\cdot 2K(\delta_*+\omega(r))(|\nabla_{\varphi}P|+\frac{C}{r}(tr_{\omega_{\varphi}}\omega_0)^{\frac{1}{2}})|\nabla_{\varphi}F|\omega_{\varphi}^n\\
&\le p(p-1)\int_Mu^p|\nabla_{\varphi}F|^2\omega_{\varphi}^n+\int_M\frac{2p^2}{p-1}K^2(\delta_*+\omega(r))^2u^p(|\nabla_{\varphi}P|^2+\frac{C'}{r^2}tr_{\omega_{\varphi}}\omega_0)\omega_{\varphi}^n.
\end{split}
\end{equation*}
For the last term in the right hand side of (\ref{4.137}):
\begin{equation*}
\begin{split}
-&\int_Mpu^{p-1}e^{K(P-\tilde{P})^2}d^cF\wedge dd^cP\wedge dP\wedge \omega_{\varphi}^{n-2}\le \int_M\frac{p}{8}u^{p-1}e^{K(P-\tilde{P})^2}|\sqrt{-1}\partial\bar{\partial}P|^2_{\varphi}\omega_{\varphi}^n\\
&+2p\int_Mu^{p-1}e^{K(P-\tilde{P})^2}|\nabla_{\varphi}F|^2|\nabla_{\varphi}P|^2\omega_{\varphi}^n.
\end{split}
\end{equation*}
Summing up,  we obtain that:
\begin{equation}\label{4.144}
\begin{split}
&\int_Mpu^{p-1}e^{K(P-\tilde{P})^2}dd^cF\wedge d^cP\wedge dP\wedge \omega_{\varphi}^{n-2}\\
&\le \int_M\frac{p(p-1)u^{p-2}}{16}|\nabla_{\varphi}u|^2\omega_{\varphi}^n+\int_M\frac{pu^{p-1}}{8}e^{K(P-\tilde{P})^2}|\sqrt{-1}\partial\bar{\partial}P|^2_{\varphi}\omega_{\varphi}^n\\
&+\int_MCp^2u^p|\nabla_{\varphi}F|^2\omega_{\varphi}^n+\int_MCpK^2(\delta_*+\omega(r))^2(|\nabla_{\varphi}P|^2+\frac{C'}{r^2}tr_{\omega_{\varphi}}\omega_0)\omega_{\varphi}^n
\end{split}
\end{equation}
Now we plug (\ref{4.142})-(\ref{4.144}) back to (\ref{4.135New}),  and we obtain:
\begin{equation}
\begin{split}
&\int_M\frac{3p(p-1)}{4}u^{p-2}|\nabla_{\varphi}u|^2\omega_{\varphi}^n\le -\int_M\frac{p}{2}u^{p-1}e^{K(P-\tilde{P})^2}|\sqrt{-1}\partial\bar{\partial}P|^2_{\varphi}\omega_{\varphi}^n\\
&-(2K-C''pK^2(\delta_*+\omega(r))^2)\int_Mpu^p|\nabla_{\varphi}P|^2\omega_{\varphi}^n+\frac{C_1(p)K}{r^2}\int_Mpu^p\omega_{\varphi}^n\\
&+\frac{\tilde{C}p}{r^2}\int_M(K(\delta_*+\omega(r))+K^2(\delta_*+\omega(r))^2)u^ptr_{\omega_{\varphi}}\omega_0\omega_{\varphi}^n+Cp\int_Mu^p|\nabla_{\varphi}F|^2\omega_{\varphi}^n.
\end{split}
\end{equation}

We are going to assume that $\delta_*$ and $r$ are sufficiently small so that:
\begin{equation}
2K-C''pK^2(\delta_*+\omega(r))^2\ge \frac{3K}{2}.
\end{equation}
Therefore we get
\begin{equation}\label{K8}
\frac{3K}{2} \int_M u^p|\nabla_{\varphi}P|^2 \omega_{\varphi}^n \le   \frac{C(p) }{r^2}\int_M u^p (1+tr_{\omega_{\varphi}}\omega_0)\omega_\varphi^n +C(p)\int_M u^p |\nabla_\varphi F|^2\omega_\varphi^n + \frac{K}{r^2}C(p)\int_M u^p \omega_\varphi^n.
\end{equation}\\
We also assume that $K(\delta_*+\omega(r))^2\le 1$,  so that:
\begin{equation*}
1\le e^{K(P-\tilde{P})^2}\le e^{K(\delta_*+\omega(r))^2}\le e.  
\end{equation*}
Therefore,  we get:
\begin{equation*}
\begin{split}
&K\int_M|\nabla_{\varphi}P|^{2(p+1)}\omega_{\varphi}^n\le \frac{C'(p)}{r^2}\int_M|\nabla_{\varphi}P|^{2p}(1+tr_{\omega_{\varphi}}\omega_0)\omega_{\varphi}^n+C'(p)\int_Mu^p|\nabla_{\varphi}F|^2\omega_{\varphi}^n\\
&+\frac{KC'(p)}{r^2}\int_Mu^p\omega_{\varphi}^n.
\end{split}
\end{equation*}

Then we just need to apply the following version of Young's inequality to the right hand side above and the claimed estimate will follow:
\begin{equation*}
a^pb\le \frac{p}{p+1}\eps a^{p+1}+\frac{1}{p+1}\eps^{-p}b^{p+1},\,a,\,b>0,\,p>0,\,\eps>0.
\end{equation*}
\end{proof}  

\end{document}
